\usepackage{framed}
\usepackage{xcolor}
\usepackage{titlesec}
\usepackage{fancyhdr}
\usepackage{microtype}
\usepackage{needspace}

% --- running heads: chapter on the left, section on the right -------------
\pagestyle{fancy}
\fancyhf{}
\renewcommand{\headrulewidth}{0.4pt}
\fancyhead[LE]{\small\itshape\nouppercase{\leftmark}}
\fancyhead[RO]{\small\itshape\nouppercase{\rightmark}}
\fancyfoot[C]{\thepage}
\fancypagestyle{plain}{\fancyhf{}\fancyfoot[C]{\thepage}\renewcommand{\headrulewidth}{0pt}}

% --- chapter openings ------------------------------------------------------
\titleformat{\chapter}[display]
  {\normalfont\sffamily\bfseries}
  {\raggedright\Large\color{black!55}\chaptertitlename\ \thechapter}
  {6pt}
  {\raggedright\huge}
\titlespacing*{\chapter}{0pt}{0pt}{28pt}
\titleformat{\section}{\normalfont\sffamily\large\bfseries}{\thesection}{0.6em}{}
\titleformat{\subsection}{\normalfont\sffamily\bfseries}{\thesubsection}{0.6em}{}

% --- keep a heading with the text under it ---------------------------------
\preto\section{\needspace{4\baselineskip}}
\preto\subsection{\needspace{3\baselineskip}}

% --- block quotes read as O'Reilly sidebars --------------------------------
\usepackage{mdframed}
\newmdenv[
  backgroundcolor=black!4, linewidth=0pt, leftmargin=0pt, rightmargin=0pt,
  innerleftmargin=10pt, innerrightmargin=10pt, innertopmargin=8pt,
  innerbottommargin=8pt, skipabove=10pt, skipbelow=10pt
]{sidebarbox}
\renewenvironment{quote}{\begin{sidebarbox}\small}{\end{sidebarbox}}

% --- characters the body font does not carry -------------------------------
% Palatino has no arrows or maths signs; these fall back to the maths font.
% Box-drawing characters are only ever inside code blocks, where Menlo has them.
\usepackage{newunicodechar}
\newunicodechar{→}{\ensuremath{\rightarrow}}
\newunicodechar{←}{\ensuremath{\leftarrow}}
\newunicodechar{≈}{\ensuremath{\approx}}
\newunicodechar{≡}{\ensuremath{\equiv}}
\newunicodechar{×}{\ensuremath{\times}}
\newunicodechar{²}{\ensuremath{^2}}
\newunicodechar{·}{\ensuremath{\cdot}}
\newunicodechar{…}{\ldots}

% --- code blocks: a tint, and air around them -------------------------------
% Two listings in a row (a program then its output) were running together into
% one grey slab; the shading and the space between are what separate them.
\definecolor{shadecolor}{gray}{0.955}
\BeforeBeginEnvironment{verbatim}{\begin{snugshade}\vspace{-2pt}}
\AfterEndEnvironment{verbatim}{\vspace{-2pt}\end{snugshade}\vspace{2pt}}
\BeforeBeginEnvironment{Shaded}{\vspace{2pt}}
\AfterEndEnvironment{Shaded}{\vspace{2pt}}
\fvset{fontsize=\small,xleftmargin=2pt}

% --- title page -------------------------------------------------------------
\usepackage{titling}
\pretitle{\vspace*{5cm}\begin{flushleft}\sffamily\bfseries\fontsize{30}{34}\selectfont}
\posttitle{\end{flushleft}\vspace{2pt}}
\preauthor{\begin{flushleft}\sffamily\large\color{black!65}}
\postauthor{\end{flushleft}}
\predate{\begin{flushleft}\sffamily\small\color{black!55}\vspace{1.5cm}}
\postdate{\end{flushleft}}

% pandoc defines shadecolor after this file is included, so the last word has
% to be said at the start of the document — otherwise highlighted listings and
% plain ones come out two different greys.
\AtBeginDocument{%
  \definecolor{shadecolor}{gray}{0.955}%
  % pandoc leaves Shaded empty unless the highlight style carries a background,
  % so highlighted listings came out unshaded next to shaded plain ones.
  \renewenvironment{Shaded}{\begin{snugshade}}{\end{snugshade}}%
}

% Inline code cannot hyphenate, so a long identifier at the end of a line used
% to push three lines a few millimetres into the margin. This lets TeX stretch
% the spaces a little further before it gives up and overflows.
\emergencystretch=3em

% --- the cover -------------------------------------------------------------
% An homage to the O'Reilly animal books, and the parts that make one
% recognisable are: a coloured band top and bottom, the engraving high on the
% page and large, and the title centred underneath it in a serif — not a
% sans-serif title hanging under a floating bird, which is what this was first.
\usepackage{graphicx}
\usepackage{tikz}
\definecolor{rillteal}{HTML}{0E7C86}
\renewcommand{\maketitle}{%
  \begin{titlepage}
  \thispagestyle{empty}
  \newgeometry{left=1.9cm,right=1.9cm,top=0cm,bottom=0cm}
  \centering
  % top band
  \noindent\begin{tikzpicture}[remember picture,overlay]
    \fill[rillteal] (current page.north west) rectangle
      ([yshift=-2.15cm]current page.north east);
    \node[anchor=west, text=white, font=\sffamily\footnotesize, xshift=1.9cm,
          yshift=-1.1cm] at (current page.north west)
      {\MakeUppercase{A functional language that compiles to small native binaries}};
  \end{tikzpicture}
  \vspace*{3.4cm}

  \includegraphics[width=0.66\textwidth]{kingfisher.png}

  \vspace{0.9cm}
  {\fontsize{33}{38}\selectfont Functional Programming\\[2pt] with Rill\par}
  \vspace{14pt}
  {\color{rillteal}\rule{0.30\textwidth}{1.2pt}}\par
  \vspace{14pt}
  {\large\itshape Writing programs out of values instead of commands\par}
  \vfill
  {\sffamily\fontsize{15}{18}\selectfont\MakeUppercase{Barış Akın}\par}
  \vspace{6pt}
  {\footnotesize\color{black!55}Rill is designed and written by Barış Akın with
   Anthropic's AI models\ \ ·\ \ barisakin@gmail.com\par}
  \vspace{10pt}
  {\footnotesize\color{black!55}Version 0.13\ \ ·\ \ every example in this book was
   compiled and run\par}
  \vspace{1.5cm}
  {\scriptsize\color{black!45}Cover: common kingfisher, from a copper plate by
   Heinrich von Kittlitz, 1832\par}
  \vspace{1.1cm}
  \noindent\begin{tikzpicture}[remember picture,overlay]
    \fill[rillteal] ([yshift=1.15cm]current page.south west) rectangle
      (current page.south east);
  \end{tikzpicture}
  \restoregeometry
  \end{titlepage}}
